%% =================================================================
%% Wei, Mei, Zhao, Xiao, Sun, Zhang, Guo.
%% "Nondestructively identifying the mechanical behavior of soft
%%  tissues using surface deformation with an explicit inverse
%%  approach."
%% Appl. Math. Modelling 134 (2024) 126--147.
%%
%% Reconstruction of page 14 (printed p. 139) as study notes.
%% Generated by: reproduce-multiple-pages-basic-tex (batch 13-16)
%%
%% Structure: Fig. 16 (Sample 3 with 5% noise) + Table 2
%% (regularization factors, 9 rows x 6 data columns) + Fig. 17
%% (bilayer ring target). No equations, no prose.
%% NOTE: Table 2 scientific notation like 5x10^-27 came from the
%% visual read — pdftotext mangled the exponents across lines.
%% =================================================================

\documentclass[11pt]{article}

\usepackage[margin=1in]{geometry}
\usepackage{amsmath, amssymb}
\usepackage{mathrsfs}
\usepackage{bm}
\usepackage{xcolor}
\usepackage{fancyhdr}
\usepackage{booktabs}
\usepackage{multirow}
\usepackage{array}

\pagestyle{fancy}
\fancyhf{}
\lhead{\small Y.~Mei et al.}
\rhead{\small Applied Mathematical Modelling 134 (2024) 126--147}
\cfoot{\small 139 $\to$ \arabic{page}}
\renewcommand{\headrulewidth}{0.4pt}

\setcounter{page}{1}

\begin{document}

% ---------- Figure 16 placeholder ----------
\noindent
\begin{center}
\fbox{\parbox{0.92\textwidth}{\centering\vspace{1em}
\textbf{[Figure 16 --- placeholder]}\\[0.8em]
Six-panel side-by-side comparison of the \textbf{Explicit Inverse
Method} (left column) vs.\ \textbf{Implicit Inverse Method} (right
column) on \textbf{Sample 3} (the three-layer target: blue bottom
$\mu^0 \approx 3$\,MPa, green middle $\mu^1 \approx 6$\,MPa, red
top $\mu^2 \approx 10$\,MPa), with \textbf{5\,\% noise} added,
using 3, 6, and 9 displacement datasets.\\[0.4em]
\textbf{Explicit (left column)}:
(a) 3 fields, SM $\in [1.00, 8.46]$\,MPa --- undershoots the
10\,MPa nominal top layer, matching the 3-field undershoot pattern
seen in Fig.~15 for Sample 2;
(b) 6 fields, SM $\in [1.00, 10.28]$\,MPa;
(c) 9 fields, SM $\in [1.00, 10.34]$\,MPa.
All three panels show wavy but distinguishable three-layer
structure.\\[0.3em]
\textbf{Implicit (right column)}:
(d) 3 fields, SM $\in [1.00, 15.49]$\,MPa --- 55\,\% overshoot,
significant green/yellow patches corrupting the top red layer;
(e) 6 fields, SM $\in [1.00, 15.98]$\,MPa --- \textbf{worse than
3 fields} on the max-SM axis, a \emph{non-monotone} degradation
that the paper's prose does not call out;
(f) 9 fields, SM $\in [1.00, 10.17]$\,MPa --- snaps back close to
nominal, cleaner interfaces.\\[0.3em]
\textbf{Takeaway}: the implicit method's colorbar max on Sample 3
under 5\,\% noise is \emph{not monotone} in the number of
displacement datasets ($15.49 \to 15.98 \to 10.17$ across 3, 6, 9
fields). This is unusual in the paper's overall story; worth
noting as a caveat to any ``more data always helps'' interpretation.
\vspace{1em}}}
\end{center}
\vspace{0.3em}
\begin{center}
\textbf{Fig. 16.} The reconstructed results with varying 3, 6 and
9 surface displacement datasets corresponding to Sample 3 (5\,\%
noise is introduced into the measured data, Unit: MPa).
\end{center}

\vspace{1em}

% ---------- Table 2 ----------
\noindent\textbf{Table 2}\\[0.2em]
{\small The values of the regularization factors utilized in
Fig.~8, Fig.~9, Fig.~10, Fig.~11, Fig.~12, Fig.~13, Fig.~14,
Fig.~15, Fig.~16.}
\vspace{0.3em}

\begin{center}
\footnotesize
\begin{tabular}{@{}ll ccc ccc@{}}
\toprule
\multicolumn{2}{l}{Inverse method}
  & \multicolumn{3}{c}{Explicit inverse method}
  & \multicolumn{3}{c}{Implicit inverse method} \\
\cmidrule(lr){3-5}\cmidrule(lr){6-8}
\multicolumn{2}{l}{Number of surface displacement datasets}
  & 3 & 6 & 9 & 3 & 6 & 9 \\
\midrule
\multirow{3}{*}{Sample 1}
  & no noise  & 0.00            & 0.00            & $5\times 10^{-27}$ & $2\times 10^{-14}$ & $5\times 10^{-14}$ & $1\times 10^{-13}$ \\
  & 1\,\% noise & $1\times 10^{-11}$ & $1\times 10^{-11}$ & $1\times 10^{-18}$ & $2\times 10^{-11}$ & $6\times 10^{-11}$ & $6\times 10^{-10}$ \\
  & 5\,\% noise & $5\times 10^{-12}$ & $1\times 10^{-10}$ & $1\times 10^{-10}$ & $1\times 10^{-10}$ & $6\times 10^{-10}$ & $6\times 10^{-9}$  \\
\midrule
\multirow{3}{*}{Sample 2}
  & no noise  & 0.00            & 0.00            & 0.00            & $1\times 10^{-14}$ & $1\times 10^{-14}$ & $1\times 10^{-14}$ \\
  & 1\,\% noise & $1\times 10^{-13}$ & $1\times 10^{-13}$ & $1\times 10^{-13}$ & $1\times 10^{-12}$ & $1\times 10^{-11}$ & $1\times 10^{-11}$ \\
  & 5\,\% noise & $1\times 10^{-10}$ & $1\times 10^{-10}$ & $1\times 10^{-10}$ & $1\times 10^{-11}$ & $1\times 10^{-10}$ & $1\times 10^{-10}$ \\
\midrule
\multirow{3}{*}{Sample 3}
  & no noise  & 0.00            & 0.00            & 0.00            & $1\times 10^{-14}$ & $1\times 10^{-14}$ & $1\times 10^{-14}$ \\
  & 1\,\% noise & 0.00            & 0.00            & 0.00            & $1\times 10^{-12}$ & $1\times 10^{-12}$ & $1\times 10^{-11}$ \\
  & 5\,\% noise & $1\times 10^{-11}$ & $1\times 10^{-11}$ & $1\times 10^{-11}$ & $1\times 10^{-11}$ & $1\times 10^{-11}$ & $1\times 10^{-10}$ \\
\bottomrule
\end{tabular}
\end{center}

\vspace{1em}

% ---------- Figure 17 placeholder ----------
\noindent
\begin{center}
\fbox{\parbox{0.80\textwidth}{\centering\vspace{1em}
\textbf{[Figure 17 --- placeholder]}\\[0.8em]
Target shear modulus distribution for the \textbf{Case 2} bilayer
ring structure: a circular ring with an outer thicker red layer
($\mu_\text{outer} = 0.386$\,MPa) surrounding an inner thinner
blue layer ($\mu_\text{inner} = 0.129$\,MPa). Black arrows point
radially inward around the perimeter, labeled ``displacement
measurement'', indicating that the inverse problem uses only the
displacement measured on the inner wall. Colorbar ticks at
$\{0.129, 0.200, 0.267, 0.333, 0.386\}$\,MPa (jet colormap,
matching the values $[37]$ for physiological soft tissue from the
paper's bibliography). The stiffness ratio outer/inner
$\approx 3.0$ is realistic for arterial wall tissue.
\vspace{1em}}}
\end{center}
\vspace{0.3em}
\begin{center}
\textbf{Fig. 17.} Target shear modulus distribution of a bilayer
ring structure (Unit: MPa).
\end{center}

% [page ends here — continues on page 15]

\end{document}
